\documentclass{beamer}

\newcommand{\week}{Banks 4 of 6}

\title{Bank Management}
\subtitle{Mishkin Chapter 9}
\author{Econ 357}
\date{\week}

% Reference the shared preamble
\input{preamble}

% Set the footer -- change 
\setbeamertemplate{footline}{
  \leavevmode%
  \vspace{2ex}
  \hbox{%
    % Left box: Econ 457
    \begin{beamercolorbox}[wd=.4\paperwidth,ht=2.5ex,dp=1ex,left]{author in head/foot}%
      \hspace{1em}Econ 357
    \end{beamercolorbox}%
    % Middle box: Week
    \begin{beamercolorbox}[wd=.2\paperwidth,ht=2.5ex,dp=1ex,center]{date in head/foot}%
      \centering\week
    \end{beamercolorbox}%
    % Right box: Slide numbers
    \begin{beamercolorbox}[wd=.4\paperwidth,ht=2.5ex,dp=1ex,center]{date in head/foot}%
      \hfill\insertframenumber{} 
    \end{beamercolorbox}%
  }%
  \vskip0pt%
}

\begin{document}

\frame{\titlepage}

\begin{frame}
    \frametitle{Outline for Banks}
        \begin{enumerate}
            \item Role of Banks
            \item Bank Balance Sheets
            \item \fbox{Bank Earnings}
            \item Bank Management (and Investment Banks)
            \item Bank Failures
            \item Bank Regulation
        \end{enumerate}
\end{frame}


\begin{frame}
    \frametitle{Outline for Today's Lecture}
    
    \begin{enumerate}
        \item Investment Banks 
        \item Bank Earnings Management
        \begin{itemize}
            \item Minimize Losses
            \item Liquidity Management
            \item Off-Balance Sheet
        \end{itemize}
        \item Three Examples: JPM, WFC, Banc of California
        \item Practice Questions
    \end{enumerate}

\end{frame}

\begin{frame}[t]
\frametitle{1. Investment Banks}
\framesubtitle{History}

    \begin{itemize}
        \item Originated in 19th century Europe and U.S. 
        \item Financed railroads, industrial expansion, and government debt
        \item Glass–Steagall Act (1933) separated commercial and investment banking 
        \begin{itemize}
            \item example: House of Morgan split into JP Morgan and Morgan Stanley
        \end{itemize}
        \item Deregulation (1990s) and globalization expanded scope and scale
    \end{itemize}
\end{frame}

\begin{frame}[t]
\frametitle{1. Investment Banks}
\framesubtitle{What Investment Banks Do}

    \begin{itemize}
        \item \textbf{Underwriting:} Issue stocks and bonds for corporations and governments
        \item \textbf{M\&A Advisory:} Advise on mergers, acquisitions, restructurings
        \item \textbf{Trading \& Market-Making:} Facilitate liquidity in financial markets
        \item \textbf{Asset Management:} Manage investments for institutions and individuals
        \item \textbf{Research:} Provide analysis of firms, industries, and markets
    \end{itemize}
\end{frame}

\begin{frame}[t]
\frametitle{1. Investment Banks}
\framesubtitle{Universal Banks}

    Today many large financial institutions operate as \textbf{universal banks}, or combined investment banks and commercial banks.   
    \vspace{0.3cm}
    
    \textbf{Commercial Banking Activities}: Accept deposits, make loans (households and firms), provide payment services
    \vspace{0.3cm}
    
    \textbf{Investment Banking Activities}: Underwriting securities, M\&A advisor, trading and market-making, asset and wealth management
    \vspace{0.3cm}

\end{frame}

\begin{frame}[t]
\frametitle{1. Investment Banks}
\framesubtitle{Universal Banks}

    There was a large amount of consolidation of investment banks and commercial banks during the Financial Crisis of 2008.   Examples include
     \begin{itemize}
            \item JPMorgan Chase acquired Bear Stearns and Washington Mutual.
            \item Bank of America acquired Merrill Lynch.
            \item Wells Fargo acquired Wachovia.
        \end{itemize}
    
    Also during 2008 some investment banks started to take deposits (example: Goldman Sachs).   
    
\end{frame}


\begin{frame}[t]
\frametitle{2. Bank Earnings Management}
\framesubtitle{Strategies}

    \textbf{Managing the bank balance sheet:}
    \begin{itemize}
        \item Liquidity Management
        \item Asset Management (credit risk, interest rate risk)
        \item Liability Management
        \item Capital Management
    \end{itemize}
    
    \textbf{Off balance sheet activities}
    \begin{itemize}
        \item Capital Markets / Trading
        \item Dealmaking
        \item Managing client assets
    \end{itemize}
\end{frame}

\begin{frame}[t]
\frametitle{2. Bank Earnings Management}
\framesubtitle{Liquidity Management}

    Banks \textbf{must} provide cash when depositors demand it.
    \vspace{1em}

    Considerations:
    \begin{itemize}
        \item Use reserves/cash (if available)
        \item Borrow an offsetting amount in public markets
        \item Borrow from the Fed (discount window)
        \item Sell securities
        \item Sell loans (most costly option, want to avoid this!)
    \end{itemize}

\end{frame}

\begin{frame}[t]
\frametitle{2. Bank Earnings Management}
\framesubtitle{Asset Management}

    Banks seek to maximize the return on their assets
    \vspace{1em}
    
    Considerations:
    \begin{itemize}
        \item Need to management credit risk
        \item Diversification can lower risk
        \item Manage liquidity
        \item Interest rate risk: asset prices change with interest rates
    \end{itemize}

\end{frame}

\begin{frame}[t]
\frametitle{2. Bank Earnings Management}
\framesubtitle{Asset Management: Credit Risk}

    Banks seek to minimize credit losses
    \vspace{1em}

    Considerations:
    \begin{itemize}
        \item Asymmetric information:
        \item Before making a loan = adverse selection
        \item After making a loan = moral hazard
        \item Tools to address – screening / underwriting and monitoring
        \item Other: long term customer relationships, collateral and compensating balances, local knowledge
    \end{itemize}

\end{frame}

\begin{frame}[t]
\frametitle{2. Bank Earnings Management}
\framesubtitle{Asset Management: Interest Rate Risk}

    Changes in interest rates causes changes in asset values and liability costs.
    \vspace{1em}


    Considerations:
    \begin{itemize}
        \item Banks frequently engage in maturity transformation: borrow short (deposits) and lend long (loans)
        \item This creates an imbalance and a vulnerability.  If interest rates increase, assets may decline in price AND the cost of their liabilities may increase.
        \item “Duration” = changes in value due to changes in interest rates.   Banks can reduce the risk by matching the "duration" of assets and liabilities.
        \item Strategies: borrow long, lend short, interest rate swaps, 
    \end{itemize}

\end{frame}

\begin{frame}[t]
\frametitle{2. Bank Earnings Management}
\framesubtitle{Liability Management}

    Banks seek to minimize the cost of funding
    \vspace{1em}
    
    Considerations:
    \begin{itemize}
        \item Deposits are low cost, but are unstable
        \item See Charles Schwab example (next slide)
        \item Match interest rate exposure
    \end{itemize}

\end{frame}

\begin{frame}[t]
\frametitle{2. Bank Earnings Management}
\framesubtitle{Liability Management}

Charles Schwab:
\begin{itemize}
    \item Had access to an inexpensive source of funding: brokerage clients left “cash” in zero yielding deposits
    \item Deposits were unstable: in March 2023, brokerage clients started to move into higher yielding CDs
    \item Cost of funding rose significantly, which in turn reduced bank profits
\end{itemize}

\centering
\includegraphics[width=0.5\textwidth]{Econ357/figures/03_banks_3_schwab.png}

\end{frame}

\begin{frame}[t]
\frametitle{2. Bank Earnings Management}
\framesubtitle{Capital Management}


    Bank management seeks to maximize return on equity.   Higher leverage can increase the return on equity (but also increase the chance of failure).
    \vspace{1em}
    
    Considerations:
    \begin{itemize}
        \item Capital ratios set by regulators (see WFC example, and more on this next lecture)
    \end{itemize}

\end{frame}

\begin{frame}[t]
\frametitle{2. Bank Earnings Management}
\framesubtitle{Capital Management}

    Wells Fargo: 
    \begin{itemize}
        \item In 2016 Wells Fargo was found to have created a number of fake accounts for its clients.
        \item The Federal Reserve punished Wells Fargo by instituting a \$1.95 tr cap on assets.
        \item The asset cap indirectly reduced leverage (WF had to turn away deposits)
        \item WFC equity value was stagnant during this period.
        \item The Federal Reserve removed the asset cap in June 2025.
        \item WFC equity increased by 30\% in 2025
    \end{itemize}

\end{frame}

\begin{frame}[t]
\frametitle{2. Bank Earnings Management}
\framesubtitle{Capital Management}

    \centering
    \includegraphics[width=0.8\textwidth]{Econ357/figures/03_banks_4_wfclev.png}

\end{frame}

\begin{frame}[t]
\frametitle{2. Bank Earnings Management}
\framesubtitle{Capital Management}

    \centering
    \includegraphics[width=\textwidth]{Econ357/figures/03_banks_4_wfcjpm.png}

\end{frame}

\begin{frame}[t]
\frametitle{2. Bank Earnings Management}
\framesubtitle{Capital Management}
    \centering
    \includegraphics[width=\textwidth]{Econ357/figures/03_banks_4_wfc2025.png}

\end{frame}

\begin{frame}[t]
\frametitle{2. Bank Earnings Management}
\framesubtitle{Off Balance Sheet Activities}

    Banks earn money in a number of ways not directly related to taking deposits and making loans.
    \vspace{1em}
    
    \begin{itemize}
        \item Investment Bank Activities (i.e. trading, dealmaking)
        \item Fee-based activities (i.e. wealth management, servicing loans)
        \item Loan originations and sales
        \item Securitization
    \end{itemize}
    \vspace{1em}

    Some of these sources of income are highly cyclical (i.e. deal making), whereas others are more stable (i.e. wealth management).   Which one do you think bank management prefers?

\end{frame}

\begin{frame}[t]
\frametitle{3. Examples}
\framesubtitle{

\includegraphics[width=0.2\textwidth]{Econ357/figures/03_banks_4_bancofca.png}}
\vspace{-1em}

           \centering
        \textbf{Net Interest Margin}
        \vspace{1em}
        

        \begin{tblr}{
        colspec = {Q[l,6cm] Q[c,2cm]}
        }
            \toprule
            Item & \% \\
            \midrule
            Average Loan Yield & 5.9\%\\
            Average Total Cost of Funds & 2.4\%\\
               Average Total Cost of Deposits & 2.1\%\\
               \midrule
            Average Net Interest Margin & 3.5\%\\
            \bottomrule
        \end{tblr}

    \blfootnote{Source: Data from 2025 Earnings Release https://investors.bancofcal.com/news-releases/news-release-details/banc-california-inc-reports-fourth-quarter-diluted-earnings-0}

\end{frame}

\begin{frame}[t]
\frametitle{3. Examples}
\framesubtitle{

\includegraphics[width=0.2\textwidth]{Econ357/figures/03_banks_4_bancofca.png}}
\vspace{-1em}


        \centering
        \textbf{Earnings}
        \vspace{1em}
        
        \begin{tblr}{
        colspec = {Q[l,6cm] Q[c,2cm]}
        }
            \toprule
            Item & \$ millions \\
            \midrule
            Total Interest Income & 1,676 \\
            Total Interest Expense &  699 \\
            Net Interest Income & 977 \\
            \midrule
            Provision for Credit Loss & 70 \\
            Noninterest Income & 142 \\
            Noninterest Expense & 735 \\
            \midrule
            Earnings Before Taxes & 313 \\
            Earnings after Taxes & 228 \\
            \bottomrule
        \end{tblr}

\end{frame}

\begin{frame}[t]
\frametitle{3. Examples}
\framesubtitle{

\includegraphics[width=0.2\textwidth]{Econ357/figures/03_banks_4_bancofca.png}}
\vspace{-1em}


        \centering
        \textbf{ROA and ROE}
        \vspace{1em}
        
        \begin{tblr}{
        colspec = {Q[l,6cm] Q[c,2cm] Q[c,2cm]}
        }
            \toprule
            Item & \$ millions & Notes\\
            \midrule
            (1) Net Earnings &228 &\\
            (2) Average Total Assets & 33,665 & \\
            (3) Return on Assets (ROA, \%) & 0.67\% & (1) / (2) \\
            (4) Average Capital (Equity) & 3,471 & \\
            (5) Capital Ratio & 10\% & (4) / (2) \\
            (6) Risk-weighted Capital Ratio &16\% &\\
            (7) Return on Equity (ROE, \%) &6.60\% &(1) / (4) \\
            \bottomrule
        \end{tblr}


\end{frame}

\begin{frame}[t]
\frametitle{3. Examples}
\framesubtitle{

\includegraphics[width=0.2\textwidth]{Econ357/figures/03_banks_4_jpm.png}}
\vspace{-1em}

       \centering
        \textbf{Earnings by Segment}
        \vspace{1em}
        
        \begin{tblr}{
        colspec = {Q[l,6cm] Q[c,2cm] Q[c,2cm]}
        }
            \toprule
            Item & \$ billion & \% of Total \\
            \midrule
            Consumer \& Community Banking & \$3,642 & 28\% \\
            Commercial \& Investment Banking & \$7,268 & 57\% \\
            Asset \& Wealth Management & \$1,808 & 14\% \\
            Total & \$12,718 & \\
            \bottomrule
        \end{tblr}

    \blfootnote{Source: Data from Q4-2025 Earnings Release https://www.jpmorganchase.com/content/dam/jpmc/jpmorgan-chase-and-co/investor-relations/documents/quarterly-earnings/2025/4th-quarter/3f2030e7-c144-4ad8-92b8-57b36851ffb6.pdf }

\end{frame}

\begin{frame}[t]
\frametitle{3. Examples}
\framesubtitle{

\includegraphics[width=0.2\textwidth]{Econ357/figures/03_banks_4_jpm.png}}
\vspace{-1em}

       \centering
        \textbf{ROA and ROE}
        \vspace{1em}
        
        \begin{tblr}{
        colspec = {Q[l,6cm] Q[c,2cm] Q[c,2cm]}
        }
            \toprule
            Item & \$ billion & Notesl \\
            \midrule
            (1) Net Earnings & 57 & \\
            (2) Average Total Assets & \$4,425 & \\
            (3) Return on Assets (ROA, \%) & 1.2\% & (1) / (2) \\
            (4) Average Capital (Equity) & 288 & \\
            (5) Capital Ratio & 6.5\% &  (4) / (2) \\
            (6) Risk-weighted Capital Ratio & 14.5\% & \tiny{Risk-weighted assets are ~\$2T} \\
            (7) Return on Equity (ROE, \%) & 17\% &(1) / (4)\\
            \bottomrule
        \end{tblr}

\end{frame}

\begin{practiceframe}[t]
\frametitle{Practice}

Regulators attempt to reduce the riskiness of banks' asset portfolios by
\begin{enumerate}[A.]
    \item limiting the amount of loans in particular categories.
    \item encouraging banks to hold risky assets such as common stocks.
    \item establishing a minimum interest rate floor that banks can earn
    \item requiring collateral for all loans.
\end{enumerate}
\vspace{1em}

For a given return on assets, the lower is bank capital
\begin{enumerate}[A.]
    \item the lower is the return for the owners of the bank.
    \item the higher is the return for the owners of the bank.
    \item the lower is the credit risk for the owners of the bank.
    \item the lower the possibility of bank failure
\end{enumerate}


\end{practiceframe}


\begin{practiceframe}[t]
\frametitle{Practice}

When a \$10 check written on the First National Bank of Chicago is deposited in an account
at Citibank, then
\begin{enumerate}[A.]
    \item the liabilities of the First National Bank decrease by \$10.
    \item the reserves of the First National Bank increase by \$10.
    \item the liabilities of Citibank decrease by \$10.
    \item the assets of Citibank decrease by \$10.
\end{enumerate}
\vspace{1em}

Which of the following are reported as liabilities on a bank's balance sheet?
\begin{enumerate}[A. ]
    \item reserves
    \item checkable deposits
    \item consumer loans
    \item deposits with other banks
\end{enumerate}

\end{practiceframe}

\begin{practiceframe}[t]
\frametitle{Practice}

Which of the following statements is FALSE?
\begin{enumerate}[A. ]
    \item A bank's assets are its uses of funds.
    \item A bank issues liabilities to acquire funds.
    \item The bank's assets provide the bank with income.
    \item Bank capital is recorded as an asset on the bank balance sheet.
\end{enumerate}

\end{practiceframe}

\end{document}